\section*{Multivibratore Triangolare Regolato}

\noindent\colorbox{orange}{
  \begin{minipage}{1.0\linewidth}
    \subsection*{Richiesta (a)}
    Descrivere qualitativamente il funzionamento del circuito, mostrando che gli stati di $V_2$ siano
    quelli di un multivibratore astabile.
  \end{minipage}
}
\subsection*{Risposta (a)}
Il circuito \'e formato da due blocchi fra cui
\begin{itemize}
\item Il primo blocco \'e un circuito \textit{sommatore invertente}.
\item Il secondo blocco \'e un circuito \textit{trigger di Schmitt}.
\end{itemize}
Il segnale prodotto dal secondo blocco \'e uno fra due stati di saturazione,
tale segnale attraversa il primo blocco, al caricarsi del condensatore la tensione
d'uscita aumenta fino ad attraversare la soglia di commutazione del trigger di Schmitt
ribaltando lo stato di saturazione (il segno), iniziando il processo di scarica del
condensatore su uno dei piatti caricando il rimanente per poi ripetere il processo. \\

\noindent\colorbox{orange}{
  \begin{minipage}{1.0\linewidth}
    \subsection*{Richiesta (b)}
    Calcolare le soglie del trigger di Schmitt che pilota l'uscita dell'astabile, ovvero i valori di
    $V_1$ in corrispondenza dei quali ha luogo la commutazione di $V_2$.
  \end{minipage}
}
\subsection*{Risposta (b)}
\begin{equation*}
  \begin{gathered}
    V_{thr+} = 3V_{OH} \\
    V_{thr-} = 3V_{OL}
  \end{gathered}
\end{equation*}

\noindent\colorbox{orange}{
  \begin{minipage}{1.0\linewidth}
    \subsection*{Richiesta (c)}
    Calcolare la pendenza del segnale $V_1$ in termini degli stati ($V_{OH}$, $V_{OL}$) di $V_2$ e del livello di
    $V_R$.
  \end{minipage}
}
\subsection*{Risposta (c)}
Si calcola la funzione di trasferimento del primo blocco. L'ingresso \textit{\underline{non}-invertente} \'e
messo a terra.
\begin{equation*}
  V_+ = 0
\end{equation*}
L'ingresso \textit{invertente}.
\begin{equation*}
  \begin{gathered}
    \frac{V_S - V_-}{Z_{R_1}} = \frac{V_- - V_R}{Z_{R_4}} + \frac{V_- - V_1}{Z_{C_1}} \\ \implies
    V_- = \frac{V_S + V_R + sRCV_1}{2 + sRC}
  \end{gathered}
\end{equation*}
nell'ipotesi che l'op-amp sia ideale e in regime di operazione lineare per il \textit{principio di corto-circuito virtuale}
\begin{equation*}
  \begin{gathered}
    V_+ = V_- \\ \implies
    \frac{V_S + V_R + sRCV_1}{2 + sRC} = 0
  \end{gathered}
\end{equation*}
la tensione in uscita dal blocco 1 \'e dunque
\begin{equation*}
  \colorbox{blu}{\boxed{V_1 = -\frac{V_2 + V_R}{sRC}}}
\end{equation*}

\noindent\colorbox{orange}{
  \begin{minipage}{1.0\linewidth}
    \subsection*{Richiesta (d)}
    In funzione di $V_R$, determinare la durata degli stati di $V_2$ assumendo per semplicit\'a che gli
    stati di saturazione degli operazionali siano uguali ed opposti (ovvero che $V_{OL} = -V_{OH}$).
  \end{minipage}
}
\subsection*{Risposta (d)}
La \textit{legge oraria} pu\'o essere utile, la si ottiene mediante la \textit{trasfromata di Laplace inversa}
\begin{equation*}
  V_1(t) = -\frac{1}{RC} \int_0^t V_2(t') + V_R(t') ~\text{d}t'
\end{equation*}
nell'ipotesi in cui i segnali siano costanti nel tempo al di fuori delle commutazioni
\begin{equation*}
  V_1(t) = -\frac{V_2 + V_R}{RC}t
\end{equation*}
da cui
\begin{equation*}
  \begin{gathered}
    -3V_{sat} = -\frac{V_{sat} + V_R}{RC}t_+ + 3V_{sat} \implies \colorbox{blu}{\boxed{t_+ = \frac{6RCV_{sat}}{V_{sat} + V_R}}} \\
    3V_{sat} = -\frac{-V_{sat} + V_R}{RC}t_- - 3V_{sat} \implies \colorbox{blu}{\boxed{t_- = \frac{6RCV_{sat}}{V_{sat} - V_R}}}
  \end{gathered}
\end{equation*}

\noindent\colorbox{orange}{
  \begin{minipage}{1.0\linewidth}
    \subsection*{Richiesta (e)}
    Ricavare la dipendenza della frequenza e del duty-cycle dello stesso segnale da $V_R$ (in
    particolare dal rapporto adimensionale $V_R/V_{OH}$) verificandone la compatibilit\'a con quanto
    misurato ai punti 2a, 2b.
  \end{minipage}
}
\subsection*{Risposta (e)}
il periodo finale
\begin{equation*}
  T = t_+ + t_- = \frac{6RCV_{sat}}{V_{sat} + V_R} + \frac{6RCV_{sat}}{V_{sat} - V_R} = \frac{12RC}{1 - \left(\frac{V_R}{V_{sat}}\right)^2}
\end{equation*}
\begin{equation*}
  \colorbox{blu}{\boxed{f = \frac{1 - \left(\frac{V_R}{V_{sat}}\right)^2}{24\pi RC}}}
\end{equation*}
il duty cycle
\begin{equation*}
  \colorbox{blu}{\boxed{D = \frac{t_+}{t_+ + t_-} = \frac{1 - \frac{V_R}{V_{sat}}}{2}}}
\end{equation*}
